\documentclass[11pt,a4paper]{report}

\usepackage[utf8]{inputenc}
\usepackage[T1]{fontenc}
\usepackage{lmodern}
\usepackage[english]{babel}
\usepackage{amsmath,amssymb,amsfonts}
\usepackage{graphicx}
\usepackage{booktabs}
\usepackage{tabularx}
\usepackage{geometry}
\usepackage{tikz}
\usepackage{listings}
\usepackage{xcolor}
\usepackage{hyperref}

\geometry{top=2.5cm, bottom=2.5cm, left=3.0cm, right=2.5cm}
\graphicspath{{../figures/}{../images/}}

\hypersetup{
    colorlinks=true,
    linkcolor=blue!70!black,
    citecolor=blue!70!black,
    urlcolor=blue!70!black
}

\lstset{
    basicstyle=\ttfamily\footnotesize,
    breaklines=true,
    frame=single,
    backgroundcolor=\color{gray!10},
    numbers=left,
    numberstyle=\tiny\color{gray}
}

\title{\textbf{Cloud-Native Network Architecture for Defense-in-Depth of Enterprise AI Services}\\
\large Mitigation of Volumetric DDoS and Semantic Resource Exhaustion}
\author{\textbf{Davide Torelli}}
\date{2026}

\begin{document}

\maketitle

\begin{abstract}
The rapid proliferation of Large Language Model (LLM) inference microservices across enterprise cloud environments has exposed modern infrastructures to critical asymmetric Denial-of-Service (DoS) and Distributed Denial-of-Service (DDoS) vulnerabilities. Due to the high computational complexity of transformer-based autoregressive token generation, an unauthenticated client transmitting sub-kilobyte HTTP requests can force hundreds of megabytes of VRAM allocation and full GPU/CPU core saturation. 

This paper presents \textbf{C-DIDA} (\textit{Cloud Defense-in-Depth for Enterprise AI}), an end-to-end, multi-tier defense architecture designed to safeguard enterprise AI endpoints against Layer 4 and Layer 7 volumetric attacks and semantic resource exhaustion. The architecture coordinates three sequential inspection tiers: (1) a perimeter reverse proxy executing early volumetric rate limiting with instantaneous sub-2ms drop latency; (2) a high-throughput API Gateway enforcing cryptographic JSON Web Token (JWT) Role-Based Access Control (RBAC) coupled with atomic Redis-backed sliding-window rate limiters; and (3) an inline heuristic Payload Inspector Intrusion Detection System (IDS) that dynamically filters repetitive token floods and oversized payloads. Experimental validation demonstrates that the architecture mitigates 95\% of hostile perimeter bursts, reducing backend server CPU saturation from 98.5\% under unconstrained attack down to 10.0\%, while preserving non-disruptive sub-4ms ingress latencies for legitimate users.
\end{abstract}

\tableofcontents
\newpage

\chapter{Introduction, Threat Modeling, and Security Objectives}
\label{chap:intro}

This chapter delineates the operational context of cloud-native enterprise Artificial Intelligence (AI) services, formalizes the threat landscape regarding computational resource exhaustion (Semantic DoS), and defines the fundamental security objectives enforced by the proposed architecture.

\section{The Landscape of Enterprise Cloud AI Applications}
Integrating Large Language Models (LLMs) and deep neural network inference into mission-critical corporate workflows has reshaped cloud service topologies. Unlike standard microservices where request latency is bounded in tens of milliseconds with nominal CPU usage, generative AI endpoints feature an acute computational asymmetry:
\begin{itemize}
    \item \textbf{Rapid VRAM and Memory Footprint}: Model parameter loading and context window allocation instantaneously allocate gigabytes of high-bandwidth memory (GPU VRAM / host RAM).
    \item \textbf{Intensive Matrix Computations}: Autoregressive token generation continuously monopolizes matrix compute units (Tensor Cores) for the entire duration of the streaming response.
    \item \textbf{KV-Cache Persistence}: Holding self-attention Key-Value (KV) cache states across concurrent user sessions binds backend system resources for extended intervals.
\end{itemize}

An unauthenticated remote adversary transmitting high-frequency sub-kilobyte requests can exhaust server computational capacity without generating significant network bandwidth, rendering traditional network-layer bandwidth filters insufficient.

\section{Threat Modeling Matrix}
To systematically categorize incoming threats, we define a multi-dimensional threat matrix targeting the ingress-to-backend continuum (Figure~\ref{fig:threat_matrix}):

\begin{figure}[htbp]
\centering
\includegraphics[width=0.95\textwidth]{figures/threat_model_mitigation.png}
\caption{Threat Modeling Matrix and Corresponding Multi-Layer Mitigation Strategy.}
\label{fig:threat_matrix}
\end{figure}

The threat categories addressed include:
\begin{enumerate}
    \item \textbf{Layer 4/Layer 7 Volumetric HTTP Flood}: High-concurrency socket floods designed to overwhelm socket descriptors and concurrency thread pools.
    \item \textbf{Credential and Quota Abuse}: Exploitation of unmetered API endpoints or unauthorized access via forged or missing tokens.
    \item \textbf{Semantic DoS via Oversized Prompts}: Prompts exceeding nominal input lengths intended to force quadratic memory expansion in self-attention mechanisms ($O(N^2)$ complexity).
    \item \textbf{Repetitive Token Spamming}: Highly redundant repetitive strings designed to invoke pathological worst-case decoding behaviors while remaining syntactically valid JSON.
    \item \textbf{Forensic Tampering}: Illegitimate alteration of audit and access logs to obscure intrusion attempts.
\end{enumerate}

\section{Security Objectives and Defense-in-Depth Principle}
The primary objective of the \textbf{C-DIDA} architecture is to enforce strict adherence to the \textit{Defense-in-Depth} paradigm (NIST SP 800-207 Zero Trust Architecture \cite{nist_800_207}). Every transaction must pass through progressively deeper inspection gates:
\begin{itemize}
    \item \textbf{Early Rejection}: Inadmissible requests must be terminated as close to the network perimeter as possible, minimizing CPU time spent on unauthenticated traffic.
    \item \textbf{Zero Unauthenticated Computation}: The inference engine must never process payloads that have not passed cryptographic signature verification and structural sanity checks.
    \item \textbf{Auditability}: Every policy violation must generate an immutable, tamper-evident forensic log.
\end{itemize}

\chapter{System Foundations, Network Hardening, and VPC Topology}
\label{chap:architecture}

This chapter details the foundational architecture of the infrastructure, spanning containerized isolation, Virtual Private Cloud (VPC) network segmentation, and TCP transport security mechanisms.

\section{Execution Environment and Virtualization}
Microservice isolation relies on Linux kernel primitives including control groups (\textit{cgroups}), network namespaces, and virtual packet filtering tables (\textit{iptables/nftables}). In our reference deployment, containerized isolation is managed via Docker Compose running over an optimized Linux kernel.

\begin{figure}[htbp]
\centering
\includegraphics[width=0.90\textwidth]{images/docker.png}
\caption{Container cluster topology running in isolated Docker environment with mapped ingress endpoints.}
\label{fig:docker_containers}
\end{figure}

As illustrated in Figure~\ref{fig:docker_containers}, the services operate within distinct container boundaries, enforcing strict memory caps and resource bounds.

\section{Virtual Private Cloud (VPC) Topology and Dual-Network Segmentation}
A cornerstone of defense-in-depth is topological decoupling between public ingress and internal data planes. As shown in Figure~\ref{fig:architettura_rete}, the architecture partitions services into two logical bridge networks:

\begin{figure}[htbp]
\centering
\includegraphics[width=0.96\textwidth]{figures/architettura_rete.png}
\caption{Cloud-Native Network Architecture and Dual-Subnet Segregation.}
\label{fig:architettura_rete}
\end{figure}

\begin{itemize}
    \item \textbf{Public DMZ Subnet (\texttt{frontend\_net})}: Exposed to external clients. It hosts the Nginx reverse proxy and the external interface of the API Gateway. Only host port \texttt{8080} is open to public traffic.
    \item \textbf{Isolated Private Subnet (\texttt{backend\_net})}: Defined with \texttt{internal: true}. Direct host routing and external internet access are entirely disabled. It accommodates the internal API Gateway interface, the AI Inference Server (\texttt{ai-application-server:5001}), and the Redis cache (\texttt{redis:6379}).
\end{itemize}

\begin{figure}[htbp]
\centering
\includegraphics[width=0.88\textwidth]{images/docker2.png}
\caption{Detailed network interface assignments across isolated bridge networks.}
\label{fig:docker_networks}
\end{figure}

\section{Transport Layer Resilience: TCP 3-Way Handshake Protection}
While IP Spoofing is prevalent in connectionless UDP-based amplification attacks, stateful TCP communication (RFC 9293 \cite{rfc9293}) enforces bidirectional sequence synchronization via the 3-Way Handshake (\texttt{SYN} $\to$ \texttt{SYN-ACK} $\to$ \texttt{ACK}). An attacker spoofing an arbitrary source IP cannot receive the server's \texttt{SYN-ACK} and is unable to finalize the handshake or dispatch Layer 7 HTTP payloads. Perimeter TCP stack controls (SYN Cookies, backlog limits) efficiently drop spoofed connection attempts at the socket level without allocating application memory.

\chapter{Three-Tier Defense-in-Depth Implementation}
\label{chap:defense_layers}

This chapter analyzes the concrete engineering implementation across the three progressive inspection layers: the perimeter reverse proxy, the role-aware API Gateway, and the inline payload inspector IDS.

\begin{figure}[htbp]
\centering
\includegraphics[width=0.95\textwidth]{figures/sequence_defense_layers.png}
\caption{Sequential Multi-Tier Inspection Workflow (Defense-in-Depth).}
\label{fig:sequence_defense_layers}
\end{figure}

As depicted in Figure~\ref{fig:sequence_defense_layers}, requests progress through sequential checkpoints. A failure at any tier triggers an immediate short-circuit termination.

\section{Tier 1: Perimeter Volumetric Rate Limiting (Nginx)}
The first line of defense is deployed on Nginx in the DMZ. Utilizing \texttt{ngx\_http\_limit\_req\_module}, it throttles incoming HTTP request volume based on client IP:

\begin{lstlisting}[language=c, caption={Nginx Ingress Configuration with Memory Rate Limit Zone}]
http {
    limit_req_zone $binary_remote_addr zone=api_limit_zone:10m rate=10r/s;

    server {
        listen 80;

        location /api/v1/ {
            limit_req zone=api_limit_zone burst=5 nodelay;
            limit_req_status 429;
            proxy_pass http://ai-api-gateway:8000/;
        }
    }
}
\end{lstlisting}

\begin{itemize}
    \item \textbf{\texttt{\$binary\_remote\_addr}}: Uses 4-byte binary IPv4 representations, tracking up to 160,000 distinct IP addresses in a compact 10MB memory zone.
    \item \textbf{\texttt{burst=5 nodelay}}: Absorbs brief legitimate spikes while instantly rejecting volumetric bursts exceeding 10 req/s with \texttt{HTTP 429} in $<2\,\text{ms}$.
\end{itemize}

\section{Tier 2: Role-Based Access Control and Atomic Sliding Window (FastAPI + Redis)}
Legitimate requests passing Tier 1 reach the asynchronous FastAPI Gateway (ASGI Uvicorn \cite{fastapi}). 

\subsection{JWT Cryptographic Authentication}
Clients authenticate via standard Bearer tokens conforming to RFC 7519 \cite{rfc7519}. The gateway decodes the claims and validates identity and user roles (\texttt{guest}, \texttt{employee}, \texttt{admin}). Unauthenticated or forged tokens trigger an immediate \texttt{HTTP 401 Unauthorized}.

\subsection{Atomic Sliding Window Algorithm on Redis ZSET}
To prevent boundary burst flaws inherent in fixed-window limiters, the Gateway enforces a dynamic sliding window backed by Redis Sorted Sets (\texttt{ZSET}) \cite{redis}:
\begin{enumerate}
    \item \textbf{Prune}: Evict timestamps older than $(t_{\text{current}} - W_{\text{window}})$ using \texttt{ZREMRANGEBYSCORE}.
    \item \textbf{Count}: Query active requests in the current window using \texttt{ZCARD}.
    \item \textbf{Admit / Drop}: If count exceeds the role quota ($Q_{\text{guest}}=5$, $Q_{\text{employee}}=60$), return \texttt{HTTP 429}; otherwise, record $t_{\text{current}}$ via \texttt{ZADD}.
\end{enumerate}
All operations are executed atomically in Redis to guarantee concurrency safety under high parallel load.

\section{Tier 3: Inline Payload Inspector IDS Heuristics}
Authenticated requests within quota limits undergo content-level inspection prior to backend dispatch:
\begin{itemize}
    \item \textbf{Input Length Cap}: Enforces a strict maximum bound ($L_{\max}=500$ characters). Prompts exceeding this cap are immediately rejected with \texttt{HTTP 400 Bad Request}, guarding against attention memory explosion.
    \item \textbf{Token Repetition Ratio}: Computes lexical vocabulary diversity:
    \begin{equation}
        \text{Repetition Ratio} = 1.0 - \frac{\text{Count}(\text{Unique Words})}{\text{Count}(\text{Total Words})}
    \end{equation}
    Prompts with a repetition ratio exceeding $0.3$ combined with high token volume trigger a semantic DoS block (\texttt{HTTP 400}).
\end{itemize}

\section{Tamper-Evident Forensic Logging (HMAC-SHA256)}
To ensure non-repudiation and forensic auditability, every security event is recorded as a structured JSON object cryptographically signed using an HMAC-SHA256 hash (RFC 2104 \cite{rfc2104}), guaranteeing tamper detection across centralized log streams.

\chapter{Experimental Validation and Performance Evaluation}
\label{chap:evaluation}

This chapter details the empirical test methodology and benchmark findings across four representative operational scenarios.

\section{Measurement Setup and Methodology}
Benchmarking was conducted by dispatching concurrent bursts of 100 requests ($N=10$ independent repetitions) using asynchronous client runners (\texttt{load-test/attack\_simulation.py} and \texttt{load-test/protected\_run.py}):
\begin{itemize}
    \item \textbf{Scenario A (Baseline Unprotected)}: 100 requests routed directly to AI Server port \texttt{5001}.
    \item \textbf{Scenario B (Guest Traffic via Perimeter)}: 100 requests sent to Nginx port \texttt{8080} with guest credentials.
    \item \textbf{Scenario C (Employee Quota)}: Evaluates employee quota limits (60 req/min) via the Gateway.
    \item \textbf{Scenario D (Admin Role)}: Evaluates privileged administrators bypassing rate limit restrictions.
\end{itemize}

\section{CPU Peak Utilization Analysis}
As illustrated in Figure~\ref{fig:benchmark_cpu}, the unprotected backend experiences immediate computational saturation, reaching a peak CPU utilization of \textbf{98.5\%}. Under Scenario B, the multi-tier defense architecture drops \textbf{95\%} of non-compliant requests at the perimeter, maintaining backend CPU utilization at an optimal \textbf{10.0\%}.

\begin{figure}[htbp]
\centering
\includegraphics[width=0.92\textwidth]{figures/benchmark_cpu_peak.png}
\caption{Backend CPU Peak Resilience and Dropped Request Ratio Across Scenarios.}
\label{fig:benchmark_cpu}
\end{figure}

\section{Latency Efficiency: Backend Processing vs. Perimeter Drop}
Figure~\ref{fig:benchmark_latency} highlights the orders-of-magnitude efficiency achieved by early rejection:
\begin{itemize}
    \item \textbf{Accepted AI Inference}: Average execution time of $\approx 2000.0\,\text{ms}$ ($\sim 2.0\,\text{s}$).
    \item \textbf{Perimeter Volumetric Drop (Nginx)}: Dropped in \textbf{1.8 ms} via \texttt{HTTP 429}.
    \item \textbf{Gateway Policy Drop (RBAC / IDS)}: Dropped in \textbf{3.5 ms} via \texttt{HTTP 429/400}.
\end{itemize}

\begin{figure}[htbp]
\centering
\includegraphics[width=0.90\textwidth]{figures/benchmark_latency_time.png}
\caption{Response Latency Comparison (Logarithmic Scale): Backend Inference vs. Tier Drops.}
\label{fig:benchmark_latency}
\end{figure}

Rejecting unauthorized traffic in under 4ms preserves computational throughput and guarantees service availability for legitimate users.

\chapter{Design Trade-offs, Limitations, and Future Work}
\label{chap:tradeoffs}

This chapter critically reviews engineering design choices, outlines current system limitations, and explores directions for future enterprise enhancements.

\section{Architectural Trade-offs}
\begin{itemize}
    \item \textbf{Redis vs. In-Memory Gateway Cache}: In-memory dictionaries avoid network round-trips but fail across distributed multi-worker Gateway deployments. Redis provides shared state and atomicity via Lua scripts at the cost of a nominal 0.5ms network hop.
    \item \textbf{Sliding Window vs. Fixed Window}: Fixed windows permit double-rate spikes at interval boundaries. Sliding-window logs prevent micro-bursts with predictable memory overhead ($O(N_{\text{active}})$).
    \item \textbf{Docker Compose vs. Kubernetes}: Docker Compose was chosen for deterministic local reproducibility and minimal networking virtualization overhead (CNI overlay), while remaining architecturally ready for Kubernetes Ingress migration.
\end{itemize}

\section{System Limitations}
\begin{enumerate}
    \item \textbf{Symmetric JWT Cryptography}: The prototype implements symmetric HMAC-SHA256 (\texttt{HS256}). Production enterprise environments should transition to asymmetric public/private key pairs (\texttt{RS256} or \texttt{ES256}) where only the identity provider holds the signing key.
    \item \textbf{Semantic Classification Breadth}: The heuristic token repetition inspector targets structural lexical spam. Detecting sophisticated adversarial prompt injections requires dedicated vector embedding guardrails (e.g., Llama Guard).
\end{enumerate}

\section{Conclusions}
The \textbf{C-DIDA} architecture demonstrates that coordinating perimeter volumetric filtering, cryptographic RBAC, and payload inspection creates a highly resilient defense for cloud-native AI microservices, dropping 95\% of hostile traffic with sub-2ms latency and slashing backend resource exhaustion by nearly 90\%.

\begin{thebibliography}{99}
\addcontentsline{toc}{chapter}{Bibliography}

\bibitem{nist_800_207}
National Institute of Standards and Technology (NIST),
\emph{Zero Trust Architecture},
NIST Special Publication 800-207, August 2020.
\url{https://doi.org/10.6028/NIST.SP.800-207}

\bibitem{owasp_llm_2023}
OWASP Foundation,
\emph{OWASP Top 10 for Large Language Model Applications (v1.1)},
Focus on \textit{LLM04: Model Denial of Service}, 2023.
\url{https://owasp.org/www-project-top-10-for-large-language-model-applications/}

\bibitem{owasp_api_2023}
OWASP Foundation,
\emph{OWASP API Security Top 10},
Focus on \textit{API4:2023 Unrestricted Resource Consumption}, 2023.
\url{https://owasp.org/www-project-api-security/}

\bibitem{rfc9293}
W.~Eddy,
\emph{Transmission Control Protocol (TCP) Specification},
RFC 9293, IETF Internet Standard, August 2022.
\url{https://datatracker.ietf.org/doc/html/rfc9293}

\bibitem{rfc8446}
E.~Rescorla,
\emph{The Transport Layer Security (TLS) Protocol Version 1.3},
RFC 8446, IETF Proposed Standard, August 2018.
\url{https://datatracker.ietf.org/doc/html/rfc8446}

\bibitem{rfc7519}
M.~Jones, J.~Bradley, and N.~Sakimura,
\emph{JSON Web Token (JWT)},
RFC 7519, IETF Proposed Standard, May 2015.
\url{https://datatracker.ietf.org/doc/html/rfc7519}

\bibitem{rfc2104}
H.~Krawczyk, M.~Bellare, and R.~Canetti,
\emph{HMAC: Keyed-Hashing for Message Authentication},
RFC 2104, IETF Informational, February 1997.
\url{https://datatracker.ietf.org/doc/html/rfc2104}

\bibitem{redis}
Redis Ltd.,
\emph{Redis Documentation: Sorted Sets (ZSET), Lua Scripting and Atomic Operations},
2024.
\url{https://redis.io/docs/}

\bibitem{fastapi}
S.~Ramírez,
\emph{FastAPI Framework Documentation: High Performance Asynchronous Web APIs with Python},
2024.
\url{https://fastapi.tiangolo.com/}

\bibitem{uvicorn}
Encode OSS,
\emph{Uvicorn: The lightning-fast ASGI server},
2024.
\url{https://www.uvicorn.org/}

\end{thebibliography}


\end{document}
